\setcounter{section}{0}
\begin{center}
	\zihao{3}
	\textbf{河北工程大学2023$\sim*$2024学年第二学期期末考试(A)卷}
\end{center}


\section{单项选择题}

\begin{enumerate}
	\item 设 $z=f(x,y)$ 在 $P_0(x_0,y_0)$ 处可微,则下列结论错误的是\underline{\hspace{2cm}}.
	
	\begin{choices}
		\item $z=f(x,y)$ 在 $P_0(x_0,y_0)$ 处连续
		\item $f_x(x_0,y_0),\,f_y(x_0,y_0)$ 存在
		\item $f_x(x,y),\,f_y(x,y)$ 在 $P_0(x_0,y_0)$ 处连续
		\item $z=f(x,y)$ 在 $P_0(x_0,y_0)$ 处沿任一射线 $l$ 方向的方向导数
		$\left.\frac{\partial f}{\partial l}\right|_{P_0}$
		都存在
	\end{choices}
	
	\item 已知区域$D=\{(x,y)\mid x^2+y^2\le a^2,\ a>0\},$
	若$\iint_D \sqrt{a^2-x^2-y^2}\,dx\,dy=\pi,$
	则 $a=$\underline{\hspace{2cm}}.
	\begin{choices}
		\item $1$
		\item $\sqrt[3]{\dfrac{3}{2}}$
		\item $\sqrt[3]{\dfrac{3}{4}}$
		\item $\sqrt[3]{\dfrac{1}{2}}$
	\end{choices}
	
	\item 若级数 $\sum_{n=1}^{\infty}u_n$ 收敛,则\underline{\hspace{2cm}}.
	\begin{choices}
		\item $\sum_{n=1}^{\infty}|u_n|$ 必收敛
		\item $\sum_{n=1}^{\infty}u_n^2$ 必收敛
		\item 部分和数列 $\{S_n\}$ 有界
		\item 以上都不正确
	\end{choices}
	
	\item 设 $L$ 为曲线 $x+2y-4=0$ 在第一象限部分,则曲线积分$\int_L e^{2x+4y+3}\,ds=$\underline{\hspace{2cm}}.
	\begin{choices}
		\item $0$
		\item $2\sqrt{5}\,e^{11}$
		\item $e^{11}$
		\item $e^4$
	\end{choices}
	
	\item 设 $f(x)$ 是周期为 $2\pi$ 的周期函数,它在 $(-\pi,\pi]$ 上的表达式为$f(x)=
		\begin{cases}
			2, & -\pi<x\le 0,\\
			x^3, & 0<x\le \pi,
		\end{cases}$
	则 $f(x)$ 的傅里叶级数在 $x=0$ 处收敛于\underline{\hspace{2cm}}.
	\begin{choices}
		\item $1$
		\item $0$
		\item $2$
		\item $-1$
	\end{choices}
\end{enumerate}

\section{填空题}

\begin{enumerate}
	\item 极限 $\displaystyle \lim_{(x,y)\to(0,0)} \frac{3-\sqrt{xy+9}}{xy}=\underline{\hspace{2cm}}$.
	
	\item 交换二次积分的积分次序:$\int_0^1 dx \int_x^1 f(x,y)\,dy\,=\underline{\hspace{2cm}}.$
	
	\item 设 $z=f(x^2+y,xy^2)$,且 $f$ 可微,则$\frac{\partial z}{\partial x}=\underline{\hspace{2cm}}.$
	
	\item 设 $L$ 是顶点为$(0,0)$,$(3,0)$,$(3,2)$的三角形正向边界,则$\oint_L(2x-y+4)\,dx+(5y+3x-6)\,dy=\underline{\hspace{2cm}}.$
	
	\item 设区域$\Omega=\{(x,y,z)\mid x^2+y^2+z^2\le 4\},$则$\iiint_{\Omega}xe^{-(x^2+y^2+z^2)}\,dx\,dy\,dz=\underline{\hspace{2cm}}.$
\end{enumerate}

\section{计算题}

\begin{enumerate}
	\item 求函数 $z=xe^{x-2y}$ 在点 $(1,\frac12)$ 处的全微分.
	
	\item 求曲面$z=x^2+y^2-1$在点 $(2,1,4)$ 处的切平面方程.
	
	\item 求函数$f(x,y)=4(x-y)-x^2-3y^2$的极值.
	
	\item 判定级数$\sum_{n=1}^{\infty}\frac{(-1)^n}{n!}$的敛散性；如果收敛,判断是绝对收敛还是条件收敛.
	
	\item 计算二重积分$\iint_D \sqrt{x^2+y^2}\,d\sigma,$其中 $D$ 是圆环型闭区域$D=\{(x,y)\mid 1\le x^2+y^2\le 4\}.$
	
	\item 设 $L$ 是抛物线 $x=y^2$ 上由点 $A(4,2)$ 到点 $B(4,-2)$ 的一段弧,计算$\int_L 2xy\,dx+x^2\,dy.$
\end{enumerate}

\section{解答题}

\begin{enumerate}

\item 求微分方程$y''-2y'-3y=3x+1$的通解.


\item 求幂级数$\sum_{n=1}^{\infty}\frac{1}{n}x^{n+1}$的收敛域及和函数.


\item 计算曲面积分$\iint_{\Sigma}\left(x^3\,dy\,dz+y^3\,dz\,dx+z^3\,dx\,dy\right),$其中 $\Sigma$ 是球面$x^2+y^2+z^2=1$的外侧.


\item 设曲线积分$I=\int_L y\big[ \varphi(x)-e^x \big] \,dx-\varphi(x)\,dy$与路径无关,其中 $\varphi(x)$ 可导且 $\varphi(0)=0$,求 $\varphi(x)$.

\end{enumerate}

\newpage